\documentclass{article}
\usepackage{amsmath, amssymb}
\usepackage{geometry}

\geometry{a4paper, margin=1in}

\begin{document}

\title{Exploration des Nombres Fondamentaux et de l'Éther-Lumière}
\author{HG}
\date{\today}

\maketitle

\section{Introduction}
Cette discussion explore les fondements mathématiques liés aux nombres récurrents tels que 360 et 129600, ainsi que leur lien potentiel avec des concepts de résonance universelle et d'éther-lumière.

\section{Formulations Mathématiques}
Les calculs initiaux posent des bases sur lesquelles une nouvelle vision peut être développée :
\begin{align*}
    u &= \frac{1}{\pi} \times 360 \times 10^8 = 1.14591559026165 \times 10^{-6} \\
    s &= \frac{1}{\pi \times 360 \times 10^8} = 8.84194128288307 \times 10^{-12} \\
    \text{\c} &= \frac{1}{\sqrt{u \cdot s}} = 314159265 \text{ rad/s} 
\end{align*}

Cette formulation permet d'imaginer une nouvelle structure physique où la vitesse limite classique disparaît, laissant place à une vision basée sur les fréquences angulaires et de résonance.

\section{Révision des Équations de Maxwell}
Une nouvelle approche des équations de Maxwell est proposée :
\begin{equation}
    \nabla E = \frac{\mu}{\epsilon}
\end{equation}
avec une valeur caractéristique :
\begin{equation}
    Z^2 = 129600
\end{equation}
Ce nombre apparaît dans plusieurs contextes et pourrait être lié à une propriété fondamentale de l’éther-lumière.

\section{Perspectives Historiques}
Des traces historiques indiquent que certaines civilisations anciennes auraient pu avoir une année de 360 jours. Une estimation basée sur le ralentissement de la rotation terrestre montre qu’une telle période aurait pu exister il y a environ 62,9 millions d’années.

\section{Conclusion}
Ces résultats et hypothèses ouvrent des pistes pour une reformulation des interactions fondamentales en physique et en cosmologie, sans recours aux paradigmes modernes du champ et de l’espace-temps.

\end{document}

